\documentclass{report}
\usepackage{amsmath,amssymb}
\setlength{\parindent}{0mm}
\setlength{\parskip}{1em}
\begin{document}
\begin{center}
\rule{6in}{1pt} \
{\large Bhuvana Srinivasan \\
{\bf Semi-implicit nonlinear solvers for the full two-fluid plasma model}}

Los Alamos National Laboratory \\ MS B268 \\ P O Box 1663 \\ Los Alamos \\ NM 87545
\\
{\tt bhuvana@lanl.gov}\\
Dana A. Knoll\end{center}

Plasma fluid models have frequently employed single-fluid descriptions
of the ion species and often neglect electron physics. While
single-fluid models have worked for a number of macroscale
applications, a number of applications in fusion as well as space
plasma physics require two-fluid descriptions of plasmas. Two-fluid
descriptions involve solving the Euler equations for each of the ion
and electron species with Maxwell's equations describing the electric
and magnetic fields. Two-fluid effects become significant when spatial
scales $L\sim \delta_i$ and temporal scales $\tau \sim 1/\omega_{ci}$
with $\delta_i$ being the ion skin depth and $\omega_{ci}$ being the
ion cyclotron frequency. Neglecting electron inertia, neglecting the
speed of light, and assuming quasineutrality leads to the more
commonly used Hall-MHD model. However, to capture local non-neutral
effects accurately, a full two-fluid plasma model is needed.

An explicit full two-fluid plasma model has been developed using the
discontinuous Galerkin method in the code WARPX (Washington
Approximate Riemann Plasma) for applications of plasma shocks,
Z-pinches, field reversed configurations, magnetic reconnection, and
plasma sheaths to name a few. The explicit implementation has severe
time-step restrictions due to the need to resolve the speed of light
and the plasma frequencies at every time-step. For simulations with
large magnetic fields, the cyclotron frequencies can be very
restrictive as well. Additionally higher spatial orders of the
explicit discontinuous Galerkin method require more restrictive CFL
conditions. Artificial ion-to-electron mass ratios and artificial
light speeds are used to obtain results with reasonable computational
effort when using an explicit two-fluid plasma model. This provides
the motivation for a semi-implicit implementation which would allow
the use of realistic parameters and geometries. The ion fluid is
solved explicitly due to the need to resolve ion physics as a minimum,
while the restrictive operators in the electron equations and
Maxwell's equations are evolved implicitly. Due to the stiff nature of
the equations, physics-based preconditioners need to be employed with
Jacobian-free Newton-Krylov (JFNK) nonlinear solvers.

A $1$-D, single-fluid, electrostatic electron fluid description is
studied first for simplicity. This system is solved using several
methods namely, using JFNK solvers, using physics-based linearized
semi-implicit solvers, and using physics-based preconditioners with
JFNK. The method of manufactured solutions is used to artificially
create a problem with large time-scale separation. Once this is
tested, it will be extended to a two-fluid description by including
explicit solutions of an ion fluid. Picard iterations are expected to
be used to couple the implicit and explicit solutions and obtain
higher order accuracy. Results will be presented using a semi-implicit
single-fluid description with the method of manufactured
solutions. Additionally, two-fluid electrostatic solutions of the ion
acoustic shock wave are expected to be presented.


\end{document}
